\chapter{Fazit}

\todo{diagramm andere gruppe einbinden}
\todo{messwerte raumparameter}



Es wurde Dosimetrie mit einem Plattenkondensator im Ionisationskammerbereich betrieben.
Dazu wurde zunächst Untersuch ab welcher Spannung das Signal nicht mehr spannungsabhängig ist. Dies entspricht dem Bereich in dem der Kondensator als Ionisationkammer fungiert. Es wurde festgestellt, dass dieser Bereich bei etwa $U_C = \SI{200}{\volt}$ beginnt. Um sicherzugehen wurde $U_C = \SI{300}{\volt}$ verwendet.
 Mit dieser Information wurde die Abhängigkeiten der emittierten Dosisleistung einer Röntgenröhre von ihren Betriebsparametern untersucht. Dabei wurde der lineare Zusammenhang zwischen Dosisleistung und Heizstrom sehr gut gezeigt. Für die Abhängigkeit von der Beschleunigungsspannung konnte festgestellt werden, dass, sobald die $K_\alpha$ Linien angeregt werden können, ein starker Anstieg in der Dosisleistung auftritt. Beides ließ sich jeweils für die Molybdän und Kupferröhre zeigen.
Weiterhin ließ sich festhalten, dass im durchstrahlten Bereich definitiv ein Sperrbereich nach Strahlenschutzverordnung eingerichtet werden müsste.\footnote{Wenn dieser Bereich bei eingeschalteter Röntgenstrahlung betretbar wäre.} 

Dann trat ein Defekt des genutzten Röntgengeräts auf, so das die weiteren Messungen, mit Ausnahme der CT-Aufnahme, zusammen mit der anderen Gruppe an einem anderen Röntgengerät durchgeführt werden mussten.

Es wurde mittels Stabdosimetern das Abstandsquadratsgesetz gut bestätigt und festgestellt, dass die gemessenen Äquivalenzdosen vergleichbar mit den durch den Plattenkondensator gemessenen sind.

Weiterhin wurde das Totzeitverhalten eines GMZ untersucht, wobei festgestellt werden konnte, dass die Totzeit sich gut durch eine nicht paralysierende Totzeit von $\tau = \SI{8.794(28)e-05}{\second}$ beschreiben lässt. Besser noch lässt sich das Totzeitverhalten durch ein hybrides, GMZ spezifisches Modell beschreiben, wobei die nicht paralysierende Komponente (vgl. \tab{deadtime_params}) etwas kleiner als bei dem reinen nicht paralysierenden Modell ist, während die paralysierende Komponente im Vergleich dazu zwei Größenordnungen kleiner liegt. Dies entspricht den Erwartungen an einen GMZ.

Im Folgenden wurde das Lambert-Beer'sche Gesetz überprüft. Hierbei traten leichte Abweichungen von der erwarteten Form der Kurve (vgl. \fig{transmission_thickness}) auf, welche möglicherweise durch ein nicht ganz einwandfrei funktionierendes Goniometer erklärt werden können, welches die exakte Positionierung der Probe stark erschwert hat. Die so bestimmen linearen Absorptionskoeffizienten von Aluminium wurden dann weiter verwendet um die Dicke eines zweiten Probensatzes zu berechenen und mit dieser Dicke die Absorptionskoeffizienten verschiedener anderer Materialien zu berechnen. Diese Absorptionskoeffizienten verhielten sich grob den Erwartungen entsprechend, auch wenn diese aufgrund der starken Energieabhängigkeit der Koeffizienten nicht genau mit Literaturwerten abgeglichen werden konnten.

Zuletzt\footnote{Nur im Protokoll. In der Durchführung fand dies vor dem Defekt des Röntgengeräts statt.} wurde ein CT-Gerät kalibriert und eine CT-Aufnahme von einem Frosch durchgeführt. Diese Aufnahme entsprach den Erwartungen.
